\chapter{Problem Setting}

We explicitly define a heatwave in order to design guidance methods. 
This means that we decide to represent a specific subset of state trajectoriesin this work.
However, our method is flexible to any kind of target: you can swap your definition of the event you want to replicate and use the guidance to generate that kind of event.
This really only affects the target definition.

This choice enables us to side step the problem that guiding towards a classifier's gradients or minimizing / maximizing a certain variable and stopping when we achieved the wanted result.
If we specify the carachteristic of the wanted result we side step this problem entirely.
The problem then shifts to retrieve a representative set of the event we have in mind.
This is however a requirement for the first methodology.
So we could say that the problem shifts towards defining an appropriate algorithm that can achieve a representative state similar to the specified set.
Making this step explicit directly enables us to then compare the results against this set,
and characterize its flaws / benefits directly.

It is really important to understand how to evaluate what we are doing here.
Since we are aiming at a specific region of interest, and we are guiding only wrt the region of interest, 
we are loosing the global coherence of the weather states.
We shall thus evaluate the ability of our method to generate locally coherent weather scenarios, over the region of interest.

In this work we focus on designing a simple guidance method and empirically test its parameters and the properties of the results generated across different settings.
We put less weight on the physical realism of the definition of the guidance target and the realism of the generated trajectories,
and defer this to a separate line of work.
Though, we have these two in mind when designing the guidance method throughout.
